\section{Teorema de finitud para los morfismos propios}
\label{section:III.3}
\label{section:III.3-fr}

\subsection{El lema de d\'evissage}
\label{subsection:III.3.1}
\label{subsection:III.3.1-fr}

\begin{definition}[3.1.1]
\label{III.3.1.1-fr}
Sea $\boldsymbol{K}$ una categoría abeliana. Se dice que un subconjunto $\boldsymbol{K}'$ del
conjunto de objetos de $\boldsymbol{K}$ es exacto si $0\in\boldsymbol{K}'$ y si, para toda
sucesión exacta
$0\to A'\to A\to A''\to0$
en $\boldsymbol{K}$ tal que dos de los objetos $A$, $A'$, $A''$ están en $\boldsymbol{K}'$, entonces el
tercero está también en $\boldsymbol{K}'$.
\end{definition}

\begin{theorem}[3.1.2]
\label{III.3.1.2-fr}
Sea $X$ un preesquema noetheriano; se designa por $\boldsymbol{K}$ la categoría abeliana
de los $\mathcal O_X$-Módulos coherentes. Sean $\boldsymbol{K}'$ un subconjunto exacto de
$\boldsymbol{K}$, $X'$ una parte cerrada del espacio subyacente a $X$. Se supone que
para toda parte cerrada irreducible $Y$ de $X'$, de punto genérico $y$,
existe un $\mathcal O_X$-Módulo $\mathcal G\in\boldsymbol{K}'$ de soporte $Y$ tal que $\mathcal G_y$ es
un $\boldsymbol{k}(y)$-espacio vectorial de dimensión $1$. Entonces todo $\mathcal O_X$-Módulo
coherente de soporte contenido en $X'$ pertenece a $\boldsymbol{K}'$ (y en particular,
si $X'=X$, se tiene $\boldsymbol{K}'=\boldsymbol{K}$).
\end{theorem}

\begin{proof}
Consideremos la propiedad siguiente $\boldsymbol{P}(Y)$ de una parte cerrada $Y$ de $X'$:
todo $\mathcal O_X$-Módulo coherente de soporte contenido en $Y$ pertenece a
$\boldsymbol{K}'$. En virtud del principio de inducción noetheriana (\textbf{0}\textsubscript{I},
2.2.2), se ve que se está reducido a demostrar que \emph{si $Y$ es una parte
cerrada de $X'$ tal que la propiedad $\boldsymbol{P}(Y')$ es verdadera para toda parte
cerrada $Y'$ de $Y$, distinta de $Y$, entonces $\boldsymbol{P}(Y)$ es verdadera}.

Sea pues $\mathcal F\in\boldsymbol{K}$ de soporte contenido en $Y$ y probemos que
$\mathcal F\in\boldsymbol{K}'$. Designemos también por $Y$ el subpreesquema cerrado reducido de
$X$ cuyo espacio subyacente es $Y$ por espacio subyacente (\textbf{I}, 5.2.1); está definido por un Ideal coherente $\mathcal J$ de $\mathcal O_X$. Se sabe (\textbf{I}, 9.3.4) que existe un
entero $n>0$ tal que $\mathcal J^n\mathcal F=0$; para $1\leq k\leq n$ se tiene
pues una sucesión exacta
\[
0\longrightarrow
\mathcal J^{k-1}\mathcal F/\mathcal J^k\mathcal F
\longrightarrow \mathcal F/\mathcal J^k\mathcal F
\longrightarrow \mathcal F/\mathcal J^{k-1}\mathcal F
\longrightarrow0
\]
de $\mathcal O_X$-Módulos coherentes (\textbf{0}\textsubscript{I}, 5.3.6 y 5.3.3);
como $\boldsymbol{K}'$ es exacto, se ve, por inducción sobre $k$, que basta mostrar
que cada uno de los
$\mathcal F_k=\mathcal J^{k-1}\mathcal F/\mathcal J^k\mathcal F$ está en
$\boldsymbol{K}'$. Se está pues reducido a probar que $\mathcal F\in\boldsymbol{K}'$ bajo la hipótesis
suplementaria de que $\mathcal J\mathcal F=0$; equivale a decir que
$\mathcal F=j_*(j^*(\mathcal F))$, donde $j$ es la inclusión canónica
$Y\to X$. Distingamos ahora dos casos:
\begin{enumerate}
\item[\textit{a)}] $Y$ es \emph{reducible}. Sean $Y=Y'\cup Y''$, siendo $Y'$ y $Y''$
partes cerradas de $Y$, distintas de $Y$; sean también $Y'$,
$Y''$ los subpreesquemas cerrados reducidos de $X$ cuyos espacios
subyacentes son $Y'$, $Y''$ respectivamente, que están definidos respectivamente por
los Ideales coherentes $\mathcal J'$, $\mathcal J''$ de $\mathcal O_X$. Pongamos
$\mathcal F'=\mathcal F\otimes_{\mathcal O_X}
(\mathcal O_X/\mathcal J')$,
$\mathcal F''=\mathcal F\otimes_{\mathcal O_X}
(\mathcal O_X/\mathcal J'')$.
Los homomorfismos canónicos $\mathcal F\to\mathcal F'$,
$\mathcal F\to\mathcal F''$ definen pues un homomorfismo
$u:\mathcal F\to\mathcal F'\oplus\mathcal F''$. Mostremos que para todo
$z\notin Y'\cap Y''$ el homomorfismo
$u_z:\mathcal F_z\to\mathcal F'_z\oplus\mathcal F''_z$ es \emph{biyectivo}.
En efecto, se tiene $\mathcal J'\cap\mathcal J''=\mathcal J$, pues la cuestión es
local y
\oldpage[III]{116}
la igualdad anterior resulta de (\textbf{I}, 5.2.1 y 1.1.5); si $z\notin Y''$, se tiene
pues $\mathcal J'_z=\mathcal J_z$, de donde $\mathcal F'_z=\mathcal F_z$ y
$\mathcal F''_z=0$, lo que establece nuestra afirmación en este caso; se razona
del mismo modo si $z\notin Y'$. Por tanto, el núcleo y el conúcleo de $u$, que están
en $\boldsymbol{K}$ (\textbf{0}\textsubscript{I}, 5.3.4), tienen sus soportes en $Y'\cap Y''$,
y están pues en $\boldsymbol{K}'$ por hipótesis; por la misma razón, $\mathcal F'$ y
$\mathcal F''$ están en $\boldsymbol{K}'$, luego también $\mathcal F'\oplus\mathcal F''$,
puesto que $\boldsymbol{K}'$ es exacto. La conclusión resulta entonces de considerar las
dos sucesiones exactas
\[
0\longrightarrow\operatorname{Im}u\longrightarrow
\mathcal F'\oplus\mathcal F''\longrightarrow
\operatorname{Coker}u\longrightarrow0
\]
\[
0\longrightarrow\operatorname{Ker}u\longrightarrow\mathcal F
\longrightarrow\operatorname{Im}u\longrightarrow0
\]
y de la hipótesis de que $\boldsymbol{K}'$ es exacto.

\item[\textit{b)}] $Y$ es irreducible, y por tanto el subpreesquema $Y$ de $X$ es
\emph{íntegro}. Si $y$ es su punto genérico, se tiene
$(\mathcal O_Y)_y=\boldsymbol{k}(y)$, y como $j^*(\mathcal F)$ es un
$\mathcal O_Y$-Módulo coherente, $\mathcal F_y=(j^*(\mathcal F))_y$ es un
$\boldsymbol{k}(y)$-espacio vectorial de dimensión finita $m$. Por hipótesis, hay un
$\mathcal O_X$-Módulo coherente $\mathcal G$ (necesariamente de soporte $Y$)
tal que $\mathcal G_y$ es un $\boldsymbol{k}(y)$-espacio vectorial de dimensión $1$. Por
tanto, hay un $\boldsymbol{k}(y)$-isomorfismo
$(\mathcal G_y)^m\xrightarrow{\sim}\mathcal F_y$, que es también un
$\mathcal O_Y$-isomorfismo, y como $\mathcal G^m$ y $\mathcal F$ son
coherentes, existe una vecindad abierta $W$ de $y$ en $X$ y un
isomorfismo $\mathcal G^m|W\xrightarrow{\sim}\mathcal F|W$
(\textbf{0}\textsubscript{I}, 5.2.7). Sea $\mathcal H$ el grafo de este isomorfismo,
que es un sub-$(\mathcal O_X|W)$-Módulo coherente de
$(\mathcal G^m\oplus\mathcal F)|W$, canónicamente isomorfo a
$\mathcal G^m|W$ y a $\mathcal F|W$; existe pues un sub-$\mathcal O_X$-%
Módulo coherente $\mathcal H_0$ de $\mathcal G^m\oplus\mathcal F$, que induce
$\mathcal H$ sobre $W$ y $0$ sobre $X\setmin Y$ puesto que $\mathcal G^m$ y
$\mathcal F$ tienen por soporte $Y$ (\textbf{I}, 9.4.7). Las restricciones
$v:\mathcal H_0\to\mathcal G^m$ y $w:\mathcal H_0\to\mathcal F$ de las
proyecciones canónicas de $\mathcal G^m\oplus\mathcal F$ son entonces
homomorfismos de $\mathcal O_X$-Módulos coherentes, que en $W$ y en
$X\setmin Y$ se reducen a isomorfismos; en otros términos, los núcleos y
conúcleos de $v$ y $w$ tienen su soporte en el conjunto cerrado
$Y\setmin(Y\cap W)$, distinto de $Y$. Pertenecen pues a $\boldsymbol{K}'$; por otra
parte, se tiene $\mathcal G^m\in\boldsymbol{K}'$ puesto que $\mathcal G\in\boldsymbol{K}'$ y $\boldsymbol{K}'$ es
exacto. Se concluye sucesivamente, por la exactitud de $\boldsymbol{K}'$, que
$\mathcal H_0\in\boldsymbol{K}'$ y después que $\mathcal F\in\boldsymbol{K}'$. C.Q.D.
\end{enumerate}
\end{proof}

\begin{corollary}[3.1.3]
\label{III.3.1.3-fr}
Supongamos que el subconjunto exacto $\boldsymbol{K}'$ de $\boldsymbol{K}$ tiene además la propiedad
de que todo factor directo coherente de un $\mathcal O_X$-Módulo coherente
$\mathcal M\in\boldsymbol{K}'$ pertenece también a $\boldsymbol{K}'$. En estas condiciones, la
conclusión de (\hyperref[III.3.1.2-fr]{3.1.2}) sigue siendo válida cuando la
condición «$\mathcal G_y$ es un espacio vectorial sobre $\boldsymbol{k}(y)$ de dimensión~$1$» se sustituye por $\mathcal G_y\neq0$ (lo que equivale a $\operatorname{Supp}(\mathcal G)=Y$).
\end{corollary}

\begin{proof}
El razonamiento del Teorema~\sref{III.3.1.2} solo debe modificarse en el caso~\textit{(b)}; ahora $\mathcal G_y$ es un espacio vectorial sobre $\boldsymbol{k}(y)$ de dimensión $q>0$ y, por tanto, se tiene un $\mathcal O_Y$-isomorfismo $(\mathcal G_y)^m\xrightarrow{\sim}(\mathcal F_y)^q$; el final
del razonamiento de (\hyperref[III.3.1.2-fr]{3.1.2}) prueba entonces que
$\mathcal F^q\in\boldsymbol{K}'$, y la hipótesis adicional sobre $\boldsymbol{K}'$ implica que
$\mathcal F\in\boldsymbol{K}'$.
\end{proof}

\subsection{El teorema de finitud: caso de los esquemas usuales}
\label{subsection:III.3.2-fr}

\begin{theorem}[3.2.1]
\label{III.3.2.1-fr}
Sean $Y$ un preesquema localmente noetheriano, $f:X\to Y$ un morfismo
propio. Para todo $\mathcal O_X$-Módulo coherente $\mathcal F$, los
$\mathcal O_Y$-Módulos $R^qf_*(\mathcal F)$ son coherentes para $q\geq0$.
\end{theorem}

\begin{proof}
Al ser la cuestión local sobre $Y$, se puede suponer $Y$ noetheriano, luego $X$
noetheriano (\textbf{I}, 6.3.7). Los $\mathcal O_X$-Módulos coherentes $\mathcal F$ para
los cuales la conclusión del teorema (\hyperref[III.3.2.1-fr]{3.2.1}) es verdadera
forman un subconjunto \emph{exacto} $\boldsymbol{K}'$ de la categoría $\boldsymbol{K}$ de los
$\mathcal O_X$-Módulos coherentes.
\oldpage[III]{117}
En efecto, sea
$0\to\mathcal F'\to\mathcal F\to\mathcal F''\to0$ una sucesión exacta de
$\mathcal O_X$-Módulos coherentes; supongamos por ejemplo que $\mathcal F'$
y $\mathcal F''$ pertenecen a $\boldsymbol{K}'$; se tiene la sucesión exacta de cohomología
\[
R^{q-1}f_*(\mathcal F'')\xrightarrow{\partial}R^qf_*(\mathcal F')
\longrightarrow R^qf_*(\mathcal F)\longrightarrow R^qf_*(\mathcal F'')
\xrightarrow{\partial}R^{q+1}f_*(\mathcal F')
\]
en la que, por hipótesis, los cuatro términos extremos son coherentes;
lo mismo ocurre pues con el término central $R^qf_*(\mathcal F)$ por
(\textbf{0}\textsubscript{I}, 5.3.4 y 5.3.3). Se muestra de la misma manera que
cuando $\mathcal F$ y $\mathcal F'$ (resp. $\mathcal F$ y
$\mathcal F''$) están en $\boldsymbol{K}'$, lo mismo ocurre con $\mathcal F''$ (resp.
$\mathcal F'$). Además, todo \emph{factor directo} coherente $\mathcal F'$
de un $\mathcal O_X$-Módulo $\mathcal F\in\boldsymbol{K}'$ pertenece también a $\boldsymbol{K}'$: en
efecto, $R^qf_*(\mathcal F')$ es entonces factor directo de
$R^qf_*(\mathcal F)$ (G, II, 4.4.4), luego es de tipo finito, y como es
cuasicoherente (\hyperref[III.1.4.10-fr]{1.4.10}), es coherente, al ser $Y$
noetheriano. En virtud de (\hyperref[III.3.1.3-fr]{3.1.3}), se está
reducido a probar que cuando $X$ es \emph{irreducible} de punto genérico
$x$, existe \emph{un} $\mathcal O_X$-Módulo coherente $\mathcal F$ que pertenece a
$\boldsymbol{K}'$, tal que $\mathcal F_x\neq0$: en efecto, si se establece este punto, se
podrá aplicar a todo subpreesquema cerrado irreducible $Y$ de $X$, pues
si $j:Y\to X$ es la inclusión canónica, $f\circ j$ es propio (\textbf{II}, 5.4.2),
y si $\mathcal G$ es un $\mathcal O_Y$-Módulo coherente de soporte $Y$,
$j_*(\mathcal G)$ es un $\mathcal O_X$-Módulo coherente tal que
\[
R^q(f\circ j)_*(\mathcal G)=R^qf_*(j_*(\mathcal G))
\]
(G, II, 4.9.1), luego se está efectivamente en las condiciones de aplicación de
(\hyperref[III.3.1.3-fr]{3.1.3}).

Ahora bien, en virtud del lema de Chow (\textbf{II}, 5.6.2), existe un preesquema
irreducible $X'$ y un morfismo \emph{proyectivo} y sobreyectivo
$g:X'\to X$ tal que $f\circ g:X'\to Y$ es \emph{proyectivo}. Existe un
$\mathcal O_{X'}$-Módulo $\mathcal L$ amplio para $g$ (\textbf{II}, 5.3.1);
apliquemos el teorema fundamental de los morfismos proyectivos
(\hyperref[III.2.2.1-fr]{2.2.1}) a $g:X'\to X$ y a $\mathcal L$:
existe pues un entero $n$ tal que
$\mathcal F=g_*(\mathcal O_{X'}(n))$ es un $\mathcal O_X$-Módulo coherente
y $R^qg_*(\mathcal O_{X'}(n))=0$ para todo $q>0$; además, como
\[
g^*(g_*(\mathcal O_{X'}(n)))\longrightarrow\mathcal O_{X'}(n)
\]
es sobreyectivo para $n$ suficientemente grande (\hyperref[III.2.2.1-fr]{2.2.1}), se
ve que se puede suponer que, en el punto genérico $x$ de $X$, se tiene
$\mathcal F_x\neq0$ (\textbf{II}, 3.4.7). Por otra parte, como $f\circ g$ es
proyectivo y $Y$ noetheriano, los
$R^p(f\circ g)_*(\mathcal O_{X'}(n))$ son \emph{coherentes}
(\hyperref[III.2.2.1-fr]{2.2.1}). Dicho esto,
$R^*(f\circ g)_*(\mathcal O_{X'}(n))$ es el resultado de una sucesión
espectral de Leray, cuyo término $E_2$ viene dado por
\[
E_2^{pq}=R^pf_*(R^qg_*(\mathcal O_{X'}(n)))~;
\]
lo anterior muestra que esta sucesión espectral es degenerada, y se sabe
entonces (\textbf{0}, 11.1.6) que $E_2^{p0}=R^pf_*(\mathcal F)$ es isomorfo a
$R^p(f\circ g)_*(\mathcal O_{X'}(n))$, lo que termina la demostración.
\end{proof}

\begin{corollary}[3.2.2]
\label{III.3.2.2}
\label{III.3.2.2-fr}
Sea $Y$ un preesquema localmente noetheriano.
Para todo morfismo propio $f:X\to Y$, la imagen directa por $f$ de cualquier $\mathcal O_X$-módulo coherente es un $\mathcal O_Y$-módulo coherente.
\end{corollary}

\begin{corollary}[3.2.3]
\label{III.3.2.3}
\label{III.3.2.3-fr}
Sea $A$ un anillo noetheriano y sea $X$ un esquema propio sobre $A$; para todo $\mathcal O_X$-módulo coherente $\mathcal F$, los $H^p(X,\mathcal F)$ son $A$-módulos de tipo finito, y existe un entero $r>0$ tal que, para todo $\mathcal O_X$-módulo coherente $\mathcal F$ y todo $p>r$, $H^p(X,\mathcal F)=0$.
\end{corollary}

\begin{proof}
La segunda afirmación ya se demostró en~\sref{III.1.4.12}; la primera se deduce del teorema de finitud~\sref{III.3.2.1}, teniendo en cuenta el Corolario~\sref{III.1.4.11}.
\end{proof}

En particular, si $X$ es un \emph{esquema algebraico propio} sobre un cuerpo $k$, entonces, para todo $\mathcal O_X$-módulo coherente $\mathcal F$, los $H^p(X,\mathcal F)$ son espacios vectoriales sobre $k$ de \emph{dimensión finita}.

\begin{corollary}[3.2.4]
\label{III.3.2.4}
\label{III.3.2.4-fr}
Sea $Y$ un preesquema localmente noetheriano y sea $f:X\to Y$ un morfismo de tipo finito.
Para todo $\mathcal O_X$-módulo coherente $\mathcal F$ cuyo soporte sea propio sobre $Y$~\sref[II]{II.5.4.10}, los $\mathcal O_Y$-módulos $R^qf_*(\mathcal F)$ son coherentes.
\end{corollary}

\oldpage[III]{118}
\begin{proof}
Como la cuestión es local sobre $Y$, se puede suponer que $Y$ es noetheriano, y lo mismo ocurre entonces con $X$~\sref[I]{I.6.3.7}.
Por hipótesis, todo subpreesquema cerrado $Z$ de $X$ cuyo espacio subyacente sea $\operatorname{Supp}(\mathcal F)$ es propio sobre $Y$; en otras palabras, si $j:Z\to X$ es la inmersión canónica, entonces $f\circ j:Z\to Y$ es propio.
Se puede suponer que $Z$ es tal que $\mathcal F=j_*(\mathcal G)$, donde $\mathcal G=j^*(\mathcal F)$ es un $\mathcal O_Z$-módulo coherente~\sref[I]{I.9.3.5}; como $R^qf_*(\mathcal F)=R^q(f\circ j)_*(\mathcal G)$ por el Corolario~\sref{III.1.3.4}, la conclusión se deduce inmediatamente del Teorema~\sref{III.3.2.1}.
\end{proof}

\subsection{Generalización del teorema de finitud (esquemas ordinarios)}
\label{subsection:III.3.3}
\label{subsection:III.3.3-fr}

\begin{proposition}[3.3.1]
\label{III.3.3.1}
\label{III.3.3.1-fr}
Sea $Y$ un preesquema noetheriano, sea $\mathcal S$ una $\mathcal O_Y$-álgebra cuasicoherente de tipo finito, graduada en grados positivos, sea $Y'=\operatorname{Proj}(\mathcal S)$, y sea $g:Y'\to Y$ el morfismo estructural.
Sea $f:X\to Y$ un morfismo propio, sea $\mathcal S'=f^*(\mathcal S)$, y sea $\mathcal M=\bigoplus_{k\in\mathbb Z}\mathcal M_k$ un $\mathcal S'$-módulo graduado cuasicoherente de tipo finito.
Entonces los \[
R^pf_*(\mathcal M)=\bigoplus_{k\in\mathbb Z}R^pf_*(\mathcal M_k)
\] son $\mathcal S$-módulos graduados de tipo finito para todo $p$.
Supongamos además que $\mathcal S$ está generada por $\mathcal S_1$; entonces, para todo $p\in\mathbb Z$, existe un entero $k_p$ tal que, para todo $k\geq k_p$ y todo $r\geq0$, se tiene
\[
R^pf_*(\mathcal M_{k+r})=\mathcal S_rR^pf_*(\mathcal M_k).
\tag{3.3.1.1}\label{III.3.3.1.1-fr}
\] 
\end{proposition}

\begin{proof}
La primera afirmación es idéntica al enunciado del Teorema~\sref{III.2.4.1}[i], en el que simplemente se sustituye «morfismo proyectivo» por «morfismo propio».
En la demostración del Teorema~\sref{III.2.4.1}[i], la hipótesis sobre $f$ solo se utilizó para demostrar (con la notación de aquella demostración) que $R^pf'_*(\widetilde{\mathcal M})$ es un $\mathcal O_{Y'}$-módulo coherente.
Bajo las hipótesis de la Proposición~\sref{III.3.3.1}, $f'$ es propio~\sref[II]{II.5.4.2}[iii], de modo que se puede retomar sin modificación la demostración del Teorema~\sref{III.2.4.1}[i], gracias al teorema de finitud~\sref{III.3.2.1}.

Para la segunda afirmación, basta observar que existe un recubrimiento abierto afín finito $(U_i)$ de $Y$ tal que las restricciones a los $U_i$ de los dos miembros de~(\hyperref[III.3.3.1.1]{3.3.1.1}) son iguales para todo $k\geq k_{p,i}$~\sref[II]{II.2.1.6}[ii]; basta tomar como $k_p$ el mayor de los $k_{p,i}$.
\end{proof}

\begin{corollary}[3.3.2]
\label{III.3.3.2}
\label{III.3.3.2-fr}
Sea $A$ un anillo noetheriano, sea $\mathfrak m$ un ideal de $A$, sea $X$ un $A$-esquema propio y sea $\mathcal F$ un $\mathcal O_X$-módulo coherente.
Entonces, para todo $p\geq0$, la suma directa $\bigoplus_{k\geq0}H^p(X,\mathfrak m^k\mathcal F)$ es un módulo de tipo finito sobre el anillo $S=\bigoplus_{k\geq0}\mathfrak m^k$; en particular, existe un entero $k_p\geq0$ tal que, para todo $k\geq k_p$ y todo $r\geq0$, se tiene
\[
H^p(X,\mathfrak m^{k+r}\mathcal F)
=\mathfrak m^rH^p(X,\mathfrak m^k\mathcal F).
\tag{3.3.2.1}\label{III.3.3.2.1-fr}
\]
\end{corollary}

\begin{proof}
Basta aplicar la Proposición~\sref{III.3.3.1} con $Y=\operatorname{Spec}(A)$, $\mathcal S=\widetilde S$ y $\mathcal M_k=\mathfrak m^k\mathcal F$, teniendo en cuenta el Corolario~\sref{III.1.4.11}.
\end{proof}

Conviene recordar que la estructura de $S$-módulo sobre $\bigoplus_{k\geq0}H^p(X,\mathfrak m^k\mathcal F)$ se obtiene considerando, para todo $a\in\mathfrak m^r$, la aplicación \[
H^p(X,\mathfrak m^k\mathcal F)
\longrightarrow H^p(X,\mathfrak m^{k+r}\mathcal F)
\] que procede del paso a cohomología de la aplicación de multiplicación $\mathfrak m^k\mathcal F\to\mathfrak m^{k+r}\mathcal F$ definida por $a$~\sref{III.2.4.1}.

\subsection{Teorema de finitud: caso de los esquemas formales}
\label{subsection:III.3.4}
\label{subsection:III.3.4-fr}

\oldpage[III]{119}
Los resultados de esta sección (salvo la definición~\sref{III.3.4.1}) no se utilizarán en el resto de este capítulo.

\begin{env}[3.4.1]
\label{III.3.4.1}
\label{III.3.4.1-fr}
Sean $\mathfrak X$ y $\mathfrak S$ dos preesquemas formales localmente noetherianos~\sref[I]{I.10.4.2}, y sea $f:\mathfrak X\to\mathfrak S$ un morfismo de preesquemas formales.
Diremos que $f$ es un morfismo \emph{propio} si
verifica las condiciones siguientes:
\begin{enumerate}
\item[$1^\circ$] \emph{$f$ es un morfismo de tipo finito} (\textbf{I}, 10.13.3).
\item[$2^\circ$] \emph{Si $\mathcal K$ es un haz de ideales de definición de $\mathfrak S$ y se pone $\mathcal J=f^*(\mathcal K)\mathcal O_{\mathfrak X}$, $X_0=(\mathfrak X,\mathcal O_{\mathfrak X}/\mathcal J)$, $S_0=(\mathfrak S,\mathcal O_{\mathfrak S}/\mathcal K)$, entonces el morfismo $f_0:X_0\to S_0$ inducido por $f$~\sref[I]{I.10.5.6} es propio}.
\end{enumerate}
Es inmediato que esta definición no depende del haz de ideales de definición $\mathcal K$ de $\mathfrak S$ considerado; en efecto, si $\mathcal K'$ es un segundo haz de ideales de definición tal que $\mathcal K'\subset\mathcal K$, y si se pone $\mathcal J'=f^*(\mathcal K')\mathcal O_{\mathfrak X}$, $X'_0=(\mathfrak X,\mathcal O_{\mathfrak X}/\mathcal J')$, $S'_0=(\mathfrak S,\mathcal O_{\mathfrak S}/\mathcal K')$, entonces el morfismo $f'_0:X'_0\to S'_0$ inducido por $f$ es tal que el diagrama
\[
\xymatrix{
X_0\ar[r]^{f_0}\ar[d]_i & S_0\ar[d]^j\\
X'_0\ar[r]_{f'_0} & S'_0
}
\]
es conmutativo, siendo $i$ y $j$ inmersiones sobreyectivas; en virtud de~\sref[II]{II.5.4.5}, es equivalente decir que $f_0$ o $f'_0$ es propio, en virtud de (\textbf{II}, 5.4.5). Se observará que, para todo $n\geq0$, si se pone $X_n=(\mathfrak X,\mathcal O_{\mathfrak X}/\mathcal J^{n+1})$, $S_n=(\mathfrak S,\mathcal O_{\mathfrak S}/\mathcal K^{n+1})$, entonces el morfismo $f_n:X_n\to S_n$ inducido por $f$~\sref[I]{I.10.5.6} es propio para todo $n$ en cuanto lo es para $n=0$~\sref[II]{II.5.4.6}.

Si $g:Y\to Z$ es un morfismo propio de preesquemas ordinarios localmente noetherianos, $Z'$ es un subconjunto cerrado de $Z$, e $Y'$ es un subconjunto cerrado de $Y$ tal que $g(Y')\subset Z'$, entonces la extensión $\widehat g:Y_{/Y'}\to Z_{/Z'}$ de $g$ a las compleciones~\sref[I]{I.10.9.1} es un morfismo propio de preesquemas formales, como se deduce de la definición y de~\sref[II]{II.5.4.5}.

Sean $\mathfrak X$ y $\mathfrak S$ dos preesquemas formales localmente noetherianos y sea $f:\mathfrak X\to\mathfrak S$ un morfismo \emph{de tipo finito} (\textbf{I}, 10.13.3); siendo las notaciones las mismas que arriba, se dice que una parte $Z$ del espacio subyacente a $\mathfrak X$ es \emph{propio} sobre $\mathfrak S$ (o propio para $f$) si, considerado como subconjunto de $X_0$, $Z$ es \emph{propio sobre $S_0$}~\sref[II]{II.5.4.10}.
Todas las propiedades de los subconjuntos propios de preesquemas ordinarios enunciadas en~\sref[II]{II.5.4.10} siguen siendo válidas para los subconjuntos propios de preesquemas formales, como se deduce inmediatamente de las definiciones.
\end{env}

\begin{theorem}[3.4.2]
\label{III.3.4.2}
\label{III.3.4.2-fr}
Sean $\mathfrak X$ y $\mathfrak Y$ preesquemas formales localmente noetherianos, y sea $f:\mathfrak X\to\mathfrak Y$ un morfismo propio.
Para todo $\mathcal O_{\mathfrak X}$-módulo coherente $\mathcal F$, los $\mathcal O_{\mathfrak Y}$-módulos $R^qf_*(\mathcal F)$ son coherentes para todo $q\geq0$.
\end{theorem}

Sea $\mathcal J$ un haz de ideales de definición de $\mathfrak Y$, sea $\mathcal K=f^*(\mathcal J)\mathcal O_{\mathfrak X}$, y consideremos los $\mathcal O_{\mathfrak X}$-módulos
\[
\mathcal F_k=\mathcal F\otimes_{\mathcal O_{\mathfrak Y}}
(\mathcal O_{\mathfrak Y}/\mathcal J^{k+1})
=\mathcal F/\mathcal K^{k+1}\mathcal F\qquad(k\geq0)
\tag{3.4.2.1}\label{III.3.4.2.1-fr}
\]
que forman evidentemente un \emph{sistema proyectivo} de $\mathcal O_{\mathfrak X}$-módulos topológicos, tal que
\oldpage[III]{120}
\[
\mathcal F=\varprojlim_k\mathcal F_k
\]
por \sref[I]{I.10.11.3}.
Por otra parte, del Teorema~\sref{III.3.4.2} se deduce que cada $R^qf_*(\mathcal F)$, por ser coherente, está dotado naturalmente de una estructura de $\mathcal O_{\mathfrak Y}$-módulo topológico~\sref[I]{I.10.11.6}, y lo mismo ocurre con los $R^qf_*(\mathcal F_k)$.
A los homomorfismos canónicos $\mathcal F\to\mathcal F_k=\mathcal F/\mathcal K^{k+1}\mathcal F$ corresponden canónicamente homomorfismos
\[
R^qf_*(\mathcal F)\longrightarrow R^qf_*(\mathcal F_k)
\]
que son necesariamente continuos para las estructuras de $\mathcal O_{\mathfrak Y}$-módulo topológico anteriores~\sref[I]{I.10.11.6} y forman un sistema proyectivo, dando en el límite un homomorfismo funtorial canónico
\[
R^qf_*(\mathcal F)\longrightarrow\varprojlim_kR^qf_*(\mathcal F_k)
\tag{3.4.2.2}\label{III.3.4.2.2-fr}
\]
que es un homomorfismo continuo de $\mathcal O_{\mathfrak Y}$-módulos topológicos.
Junto con el Teorema~\sref{III.3.4.2}, demostraremos el
\begin{corollary}[3.4.3]
\label{III.3.4.3}
\label{III.3.4.3-fr}
Cada uno de los homomorfismos~(\hyperref[III.3.4.2.2]{3.4.2.2}) es un isomorfismo topológico.
Además, si $\mathfrak Y$ es noetheriano, el sistema proyectivo $(R^qf_*(\mathcal F/\mathcal K^{k+1}\mathcal F))_{k\geq0}$ satisface la condición \emph{(ML)}~\sref[0]{0.13.1.1}.
\end{corollary}
Comenzaremos estableciendo el Teorema~\sref{III.3.4.2} y el Corolario~\sref{III.3.4.3} cuando $\mathfrak Y$ es un esquema formal afín noetheriano~\sref[I]{I.10.4.1}:
\begin{corollary}[3.4.4]
\label{III.3.4.4}
\label{III.3.4.4-fr}
Bajo las hipótesis del Teorema~\sref{III.3.4.2}, supongamos además que $\mathfrak Y=\operatorname{Spf}(A)$, donde $A$ es un anillo ádico noetheriano.
Sea $\mathfrak J$ un ideal de definición de $A$, y pongamos $\mathcal F_k=\mathcal F/\mathfrak J^{k+1}\mathcal F$ para $k\geq0$.
Entonces los $H^n(\mathfrak X,\mathcal F)$ son $A$-módulos de tipo finito; el sistema proyectivo $(H^n(\mathfrak X,\mathcal F_k))_{k\geq0}$ satisface la condición \emph{(ML)} para todo $n$; si se pone
\[
N_{n,k}=\operatorname{Ker}\bigl(H^n(\mathfrak X,\mathcal F)
\longrightarrow H^n(\mathfrak X,\mathcal F_k)\bigr)
\tag{3.4.4.1}\label{III.3.4.4.1-fr}
\]
(igual también a $\operatorname{Im}(H^n(\mathfrak X,\mathfrak J^{k+1}\mathcal F)
\to H^n(\mathfrak X,\mathcal F))$ por la sucesión exacta de cohomología), entonces los $N_{n,k}$ definen sobre $H^n(\mathfrak X,\mathcal F)$ una filtración $\mathfrak J$-buena~\sref[0]{0.13.7.7}; finalmente, el homomorfismo canónico
\[
H^n(\mathfrak X,\mathcal F)
\longrightarrow\varprojlim_kH^n(\mathfrak X,\mathcal F_k)
\tag{3.4.4.2}\label{III.3.4.4.2-fr}
\]
es un isomorfismo topológico para todo $n$ (el miembro izquierdo está dotado de la topología $\mathfrak J$-ádica y los $H^n(\mathfrak X,\mathcal F_k)$ de la topología discreta).
\end{corollary}

Pongamos
\[
S=\operatorname{gr}(A)=\bigoplus_{k\geq0}
\mathfrak J^k/\mathfrak J^{k+1},\qquad
\mathcal M=\operatorname{gr}(\mathcal F)=\bigoplus_{k\geq0}
\mathfrak J^k\mathcal F/\mathfrak J^{k+1}\mathcal F.
\tag{3.4.4.3}\label{III.3.4.4.3-fr}
\]
Se sabe que $\mathfrak J^\Delta$ es un haz de ideales de definición de $\mathfrak Y$~\sref[I]{I.10.3.1}; sean $\mathcal K=f^*(\mathfrak J^\Delta)\mathcal O_{\mathfrak X}$, $X_0=(\mathfrak X,\mathcal O_{\mathfrak X}/\mathcal K)$, $Y_0=(\mathfrak Y,\mathcal O_{\mathfrak Y}/\mathfrak J^\Delta)
=\operatorname{Spec}(A_0)$, con $A_0=A/\mathfrak J$.
Está claro que los $\mathcal M_k=\mathfrak J^k\mathcal F/\mathfrak J^{k+1}\mathcal F$ son $\mathcal O_{X_0}$-módulos coherentes~\sref[I]{I.10.11.3}.
Consideremos, por otra parte, la $\mathcal O_{X_0}$-álgebra graduada cuasicoherente
\[
\mathcal S=\mathcal O_{X_0}\otimes_{A_0}S.
\tag{3.4.4.4}\label{III.3.4.4.4-fr}
\]
La multiplicación define un homomorfismo sobreyectivo de álgebras graduadas
\[
\mathcal S\longrightarrow
\operatorname{gr}(\mathcal O_{\mathfrak X})
=\bigoplus_{k\geq0}\mathcal K^k/\mathcal K^{k+1},
\]
por medio del cual $\mathcal M$ es un $\mathcal S$-módulo graduado.

\begingroup
\microtypesetup{protrusion=false,expansion=false}
La hipótesis de que $\mathcal F$ es un $\mathcal O_{\mathfrak X}$-módulo de tipo finito implica primero que $\mathcal M$ es
\oldpage[III]{121}
un $\mathcal S$-módulo graduado \emph{de tipo finito}.
En efecto, la cuestión es local sobre $\mathfrak X$, y se puede pues suponer, para tratarla, que $\mathfrak X=\operatorname{Spf}(B)$, donde $B$ es un anillo ádico noetheriano, y que $\mathcal F=N^\Delta$, donde $N$ es un $B$-módulo de tipo finito~\sref[I]{I.10.10.5}.
Pongamos $B_0=B/\mathfrak J B$ y $N_0=N/\mathfrak JN$.
Entonces $X_0=\operatorname{Spec}(B_0)$, y los $\mathcal O_{X_0}$-módulos cuasicoherentes $\mathcal S$ y $\mathcal M$ están asociados respectivamente a
\[
S'=\bigoplus_{k\geq0}
((\mathfrak J^k/\mathfrak J^{k+1})\otimes_{A_0}B_0)\quad\text{y}\quad
M'=\bigoplus_{k\geq0}\mathfrak J^kN/\mathfrak J^{k+1}N.
\]
Para todo $k\geq0$, la multiplicación induce una aplicación sobreyectiva
\[
(\mathfrak J^k/\mathfrak J^{k+1})\otimes_{A_0}N_0
\longrightarrow\mathfrak J^kN/\mathfrak J^{k+1}N,
\qquad \overline a\otimes\overline n\longmapsto\overline{an}.
\]
Estas aplicaciones definen un homomorfismo sobreyectivo de $S'$-módulos graduados $S'\otimes_{B_0}N_0\longrightarrow M'$.
Como $N_0$ es un $B_0$-módulo de tipo finito, $M'$ es, por tanto, un $S'$-módulo de tipo finito, lo que prueba nuestra afirmación~\sref[I]{I.1.3.13}.
\par\endgroup

Como el morfismo $f_0:X_0\to Y_0$ es \emph{propio} por hipótesis, se puede aplicar la Proposición~\sref{III.3.3.1} a $\mathcal S$, $\mathcal M$ y al morfismo $f_0$: teniendo en cuenta el Corolario~\sref{III.1.4.11}, se concluye que, para \emph{todo $n\geq0$}, $\bigoplus_{k\geq0}H^n(X_0,\mathcal M_k)$ es un $S$-módulo graduado \emph{de
tipo finito}. Esto prueba que la condición $(\mathrm F_n)$ de (\textbf{0}, 13.7.7) se
verifica para \emph{todo $n\geq0$}, cuando se considera el sistema proyectivo estricto $(\mathcal F/\mathfrak J^k\mathcal F)_{k\geq0}$ de haces de grupos abelianos sobre $X_0$, cada uno dotado de su estructura natural de «$A$-módulo filtrado». Se puede pues aplicar (\textbf{0}, 13.7.7), que prueba que:
\begin{enumerate}
\item[$1^\circ$] El sistema proyectivo
$(H^n(\mathfrak X,\mathcal F_k))_{k\geq0}$ verifica la condición (ML).
\item[$2^\circ$] Si
$H'^{n}=\varprojlim_kH^n(\mathfrak X,\mathcal F_k)$, entonces $H'^{n}$ es un $A$-módulo de tipo finito.
\item[$3^\circ$] La filtración definida sobre $H'^{n}$ por los núcleos de los homomorfismos canónicos $H'^{n}\to H^n(\mathfrak X,\mathcal F_k)$ es $\mathfrak J$-buena.
\end{enumerate}

Observemos, por otra parte, que si se pone $X_k=(\mathfrak X,\mathcal O_{\mathfrak X}/\mathcal K^{k+1})$, entonces $\mathcal F_k$ es un $\mathcal O_{X_k}$-módulo coherente~\sref[I]{I.10.11.3}; y si $U$ es un abierto afín de $X_0$, entonces $U$ es también un abierto afín de cada $X_k$~\sref[I]{I.5.1.9}, de modo que $H^n(U,\mathcal F_k)=0$ para todo $n>0$ y todo $k$~\sref{III.1.3.1}, y $H^0(U,\mathcal F_k)\to H^0(U,\mathcal F_h)$ es sobreyectivo para $h\leq k$~\sref[I]{I.1.3.9}.
Se está, pues, en las condiciones de~\sref[0]{0.13.3.2}, y la aplicación de~\sref[0]{0.13.3.1} demuestra que $H'^{n}$ se identifica canónicamente con $H^n(\mathfrak X,\varprojlim_k\mathcal F_k)=H^n(\mathfrak X,\mathcal F)$; esto termina la demostración del Corolario~\sref{III.3.4.4}.

\begin{env}[3.4.5]
\label{III.3.4.5}
\label{III.3.4.5-fr}
Volvamos a la demostración de~\sref{III.3.4.2} y~\sref{III.3.4.3}.
Demostremos primero las proposiciones en el caso $\mathfrak Y=\operatorname{Spf}(A)$ considerado en~\sref{III.3.4.4}; para ello, para todo $g\in A$, apliquemos~\sref{III.3.4.4} al esquema formal afín noetheriano inducido sobre el abierto $\mathfrak Y_g=\mathfrak D(g)$ de $\mathfrak Y$, que es igual a $\operatorname{Spf}(A_{\{g\}})$, y al preesquema formal inducido por $\mathfrak X$ sobre $f^{-1}(\mathfrak Y_g)$; observemos que $\mathfrak Y_g$ es también un abierto afín del preesquema $Y_k=(\mathfrak Y,\mathcal O_{\mathfrak Y}/(\mathfrak J^\Delta)^{k+1})$, y como $\mathcal F_k$ es un $\mathcal O_{X_k}$-módulo coherente, se tiene
\[
H^n(f^{-1}(\mathfrak Y_g),\mathcal F_k)
=\Gamma(\mathfrak Y_g,R^nf_*(\mathcal F_k))
\]
para todo $k\geq0$, en virtud del Corolario~\sref{III.1.4.11}.
El homomorfismo canónico
\[
H^n(f^{-1}(\mathfrak Y_g),\mathcal F)
\longrightarrow\varprojlim_k\Gamma(\mathfrak Y_g,R^nf_*(\mathcal F_k))
\]
es un isomorfismo; pero se tiene~\sref[0]{0.3.2.6}
\[
\varprojlim_k\Gamma(\mathfrak Y_g,R^nf_*(\mathcal F_k))
=\Gamma(\mathfrak Y_g,\varprojlim_kR^nf_*(\mathcal F_k))
\]
\oldpage[III]{122}
y como el haz $R^nf_*(\mathcal F)$ es el haz asociado al prehaz $\mathfrak Y_g\leadsto H^n(f^{-1}(\mathfrak Y_g),\mathcal F)$ sobre los $\mathfrak Y_g$~\sref[0]{0.3.2.1}, se ha demostrado que el homomorfismo~(\hyperref[III.3.4.2.2]{3.4.2.2}) es \emph{biyectivo}.
Probemos después que $R^nf_*(\mathcal F)$ es un $\mathcal O_{\mathfrak Y}$-módulo coherente y, más precisamente, que se tiene
\[
R^nf_*(\mathcal F)=(H^n(\mathfrak X,\mathcal F))^\Delta.
\tag{3.4.5.1}\label{III.3.4.5.1-fr}
\]

Con la notación anterior, como $\mathcal F_k$ es un $\mathcal O_{X_k}$-módulo coherente~\sref{III.1.4.13}, se tiene
\[
\Gamma(\mathfrak Y_g,R^nf_*(\mathcal F_k))
=(\Gamma(\mathfrak Y,R^nf_*(\mathcal F_k)))_g
=(H^n(\mathfrak X,\mathcal F_k))_g.
\]

Ahora bien, los $H^n(\mathfrak X,\mathcal F_k)$ forman un sistema proyectivo que satisface (ML), y su límite proyectivo $H^n(\mathfrak X,\mathcal F)$ es un $A$-módulo de tipo finito.
Se concluye~\sref[0]{0.13.7.8} que se tiene
\[
\varprojlim_k((H^n(\mathfrak X,\mathcal F_k))_g)
=H^n(\mathfrak X,\mathcal F)\otimes_AA_{\{g\}}
=\Gamma(\mathfrak Y_g,(H^n(\mathfrak X,\mathcal F))^\Delta)
\]
teniendo en cuenta~\sref[I]{I.10.10.8} aplicado a $A$ y $A_{\{g\}}$; esto demuestra~(\hyperref[III.3.4.5.1]{3.4.5.1}), pues $\Gamma(\mathfrak Y_g,R^nf_*(\mathcal F))
=\varprojlim_k\Gamma(\mathfrak Y_g,R^nf_*(\mathcal F_k))$.

Como (\hyperref[III.3.4.2.2]{3.4.2.2}) es entonces un isomorfismo de $\mathcal O_{\mathfrak Y}$-módulos coherentes, es necesariamente un isomorfismo \emph{topológico}~\sref[I]{I.10.11.6}.
Finalmente, de las relaciones $R^nf_*(\mathcal F_k)=(H^n(\mathfrak X,\mathcal F_k))^\Delta$ se deduce que el sistema proyectivo $(R^nf_*(\mathcal F_k))_{k\geq0}$ satisface (ML)~\sref[I]{I.10.10.2}.

Una vez demostrados \sref{III.3.4.2} y \sref{III.3.4.3} en el caso en que el preesquema formal $\mathfrak Y$ es afín noetheriano, el paso al caso general es inmediato para~\sref{III.3.4.2} y para la primera afirmación de~\sref{III.3.4.3}, que son locales sobre $\mathfrak Y$.
En cuanto a la segunda afirmación de~\sref{III.3.4.3}, como $\mathfrak Y$ es noetheriano, basta recubrirlo por un número finito de abiertos afines noetherianos $U_i$ y observar que las restricciones del sistema proyectivo $(R^qf_*(\mathcal F_k))$ a cada uno de los $U_i$ satisfacen (ML).
\end{env}

En el curso de la demostración hemos probado además:
\begin{corollary}[3.4.6]
\label{III.3.4.6}
\label{III.3.4.6-fr}
Bajo las hipótesis del Corolario~\sref{III.3.4.4}, el homomorfismo canónico
\[
H^q(\mathfrak X,\mathcal F)
\longrightarrow\Gamma(\mathfrak Y,R^qf_*(\mathcal F))
\tag{3.4.6.1}\label{III.3.4.6.1-fr}
\]
es biyectivo.
\end{corollary}
